\definecolor{mygreen}{rgb}{0,0.6,0}
\definecolor{mygray}{rgb}{0.5,0.5,0.5}
\definecolor{mypurple}{rgb}{0.58,0,0.82}

\lstset{
  frame=single,
  numbers=none,
  breaklines=true,
  captionpos=b,
  basicstyle=\ttfamily\footnotesize,
  commentstyle=\color{mygreen},
  keywordstyle=\color{blue},
  stringstyle=\color{mypurple},
}

\section{Išeities kodas ir paleidimas}
\label{chap:priedas-paleidimas}

Emuliatoriaus išeities kodas pasiekiamas Github aplinkoje -- \href{https://github.com/samariezas/remu/tree/qpu}{https://github.com/samariezas/remu/tree/qpu}
Surinktas emuliatorius gali būti parsisiųstas ir paleistas Linux aplinkoje naudojant šias komandas:

\begin{lstlisting}[language=Bash]
wget https://mif.vu.lt/~jope9155/remu/remu.AppImage
wget https://mif.vu.lt/~jope9155/remu/remu-package.tar.gz
chmod +x remu.AppImage
tar xvf remu-package.tar.gz
remu-package/run.sh
\end{lstlisting}

Emuliatorius sukompiliuotas Sandy Bridge procesorių architektūrai norint užtikrinti korektišką veikimą VU MIF Linux aplinkoje. Emuliatorius (bei vaizdas: Linux, initramfs, OpenSBI ir reikalingi kryžminiai kompiliatoriai) gali būti sukompiliuotas naudojant Nix:

\begin{lstlisting}
nix run github:samariezas/remu/qpu
\end{lstlisting}

Paleistoje sistemoje pateiktos dvi pavyzdinės programos, kurias galima sukompiliuoti su \texttt{tcc}. Naujas programas galima rašyti naudojant \texttt{vi} programą. Programos kurios naudoja QPU turi būti susietos naudojant \texttt{-lqpu}:

\begin{lstlisting}
/samples # ls
build-all.sh  hello.c       qpu-basic.c
/samples # ./build-all.sh 
+ tcc hello.c -o hello
+ tcc qpu-basic.c -lqpu -o qpu-basic
/samples # ./hello
Hello, World!
/samples # ./qpu-basic
QPU register id: 1
q0=0, q1=1, q2=1
/samples # ./qpu-basic
QPU register id: 1
q0=0, q1=1, q2=1
/samples # ./qpu-basic
QPU register id: 1
q0=1, q1=0, q2=0
/samples #
\end{lstlisting}

\begin{figure}[h]
    \centering
\begin{quantikz}[row sep=0.5cm, column sep=0.55cm]
\lstick{$\ket{0}_{q_0}$}
  & \gate{H}
  & \ctrl{1}
  & \ctrl{1}
  & \qw
  & \qw
  & \gate{Z}
  & \meter{} \cw \rstick{$r_0$}
\\
\lstick{$\ket{0}_{q_1}$}
  & \qw
  & \targ{}
  & \ctrl{1}
  & \qw
  & \gate{Y}
  & \qw
  & \meter{} \cw \rstick{$r_1$}
\\
\lstick{$\ket{0}_{q_2}$}
  & \qw
  & \qw
  & \targ{}
  & \gate{X}
  & \qw
  & \qw
  & \meter{} \cw \rstick{$r_2$}
\end{quantikz}
    \caption{Programos {\normalfont\ttfamily qpu-basic.c} kvantinės grandinės atvaizdavimas.}
    \label{fig:qpu-basic-diagrama}
\end{figure}


Žemiau pateiktame kodo fragmente matoma pavyzdinė QPU programa \texttt{qpu-basic.c}, kurios kvantinės grandinės schema parodyta \ref{fig:qpu-basic-diagrama} paveiksliuke.
\begin{lstlisting}[language=C]
#include <stdio.h>
#include <qpu.h>

int main(){
    uint64_t id = qpu_new_qureg(0, 3);
    printf("QPU register id: %llu\n", id);
    qpu_hadamard(0, id);
    qpu_cnot(0, 1, id);
    qpu_toffoli(0, 1, 2, id);

    qpu_sigma_x(2, id);
    qpu_sigma_y(1, id);
    qpu_sigma_z(0, id);

    int r2 = qpu_bmeasure(2, id);
    int r1 = qpu_bmeasure(1, id);
    int r0 = qpu_bmeasure(0, id);
    printf("q0=%d, q1=%d, q2=%d\n", r0, r1, r2);
    return 0;
}
\end{lstlisting}
